% GENERATED FILE. Edit canonical lessons, then run npm run generate:books.
% canonical-source-hash: fnv1a64:96a2ebf8a428c459
% canonical-lessons: ZH-C226-caodi, ZH-W226-cao-grass, ZH-C226-hu, ZH-C226-senlin, ZH-C226-caoyuan, ZH-C226-fengjing

\chapter{Grass and Lakes}
\label{ch:zh-grass-and-lakes}
\hlchaptermodality{226}

\begin{chapteropening}
\textbf{By the end of this chapter:} \emph{I can describe a lawn, a lake, a forest, grassland and the scenery.}

Complete the last lesson of chapter 226: I can describe a lawn, a lake, a forest, grassland and the scenery.
\end{chapteropening}

\section[\texorpdfstring{\zh{草地} (cǎodì)}{草地 (cǎodì)}]{\zh{草地} (cǎodì) — lawn; grassy ground}

\label{lesson:ZH-C226-caodi}



\begin{quote}
Before the new one: say the Mandarin for to pick fruit, then the Mandarin for to camp out.
\end{quote}

\subsection*{You'll want to know: \zh{草地}}
\textbf{\zh{草地}} — \emph{cǎodì} — \textquotedblleft{}lawn; grassy ground\textquotedblright{}.

\begin{quote}
\emph{gōngyuán lǐ yǒu yí kuài dà cǎodì} — There's a big lawn in the park.
\end{quote}

The new character, in pieces: \zh{草} — nine strokes: the grass piece of \zh{菜} on top, then \zh{早} under it. The next lesson puts it in your hand.

\begin{cousinweb}
\textbf{\zh{草地}} (\emph{cǎodì}) is built from \zh{草} 'grass' and \zh{地} 'ground', both inherited from Middle Chinese, \emph{tshawX} and \emph{dijH} (Guangyun, in Baxter's transcription), which glosses the first 'all the plants'.
\end{cousinweb}

\begin{grammarlens}[title={What you've built}]
One more word for describing things.
\end{grammarlens}

\subsection*{Your turn}
\begin{itemize}
\raggedright
  \item \emph{Say it:} \emph{cǎodì}
  \item \emph{Say it:} \emph{cǎodì}, once more
  \item \emph{Say it:} \emph{gōngyuán lǐ yǒu yí kuài dà cǎodì} — There's a big lawn in the park.
  \item \emph{Recall:} say the Mandarin for to pick fruit, then the Mandarin for to camp out, then say \emph{cǎodì} again
\end{itemize}

\begin{tcolorbox}[breakable,colback=teal!4,colframe=teal!35!black,title={Before you move on}]
What does \zh{草地} mean? (\textquotedblleft{}Lawn; grassy ground\textquotedblright{}.) Say it once more.
\end{tcolorbox}
\section[\texorpdfstring{\zh{草} (cǎo)}{草 (cǎo)}]{\zh{草} — the grass piece over \zh{早}}

\label{lesson:ZH-W226-cao-grass}



\begin{quote}
The word first, then the character inside it.

\begin{itemize}
\raggedright
  \item \emph{Recall:} say \emph{cǎodì} — \textbf{R1}, one lesson back, before it decays
\end{itemize}

\begin{quote}
\zh{早} \zh{花}
\end{quote}

\emph{Read it:} each one aloud, with nothing said first — \emph{zǎo}, \emph{huā}
\end{quote}

\subsection*{Script — the grass piece over \zh{早}}
\hlblockfigure{\detokenize{figures/ZH-W226-cao-grass-filmstrip.pdf}}{How \zh{草} is written, stroke by stroke}

\begin{quote}
\zh{草}
\end{quote}

\textbf{Two pieces, top before bottom.}

\textbf{The grass piece, three strokes,} as in \zh{菜} and \zh{花}: a bar across, then two short verticals.

\textbf{Under it, \zh{早}, six strokes,} as in \zh{早上}: \zh{日}, four strokes, then a long bar across and a vertical down through its middle.

3 + 6 = nine.

\subsection*{Your turn}
\begin{itemize}
\raggedright
  \item \emph{Write it:} \zh{草}, all nine, with the model in front of you
  \item \emph{Cover:} write it again from memory, and count the pieces, not the strokes
\end{itemize}

\begin{tcolorbox}[breakable,colback=teal!4,colframe=teal!35!black,title={Before you move on}]
How many strokes in \zh{草}? (\textbf{Nine}.) What sits under the grass piece? (\textbf{\zh{早}}.) What comes last? (\textbf{The long vertical through the middle}.)

Source: \href{https://raw.githubusercontent.com/chanind/hanzi-writer-data/68d10a4b21150cae5e1ebbd223eed289cf32d90c/data/\%E8\%8D\%89.json}{Hanzi Writer Data \zh{草}, pinned stroke paths and medians}.
\end{tcolorbox}
\section[hú]{hú — lake}

\label{lesson:ZH-C226-hu}



\begin{quote}
Before the new one: say the Mandarin for to camp out, then the Mandarin for a lawn.
\end{quote}

\subsection*{You'll want to know: hú}
\emph{hú} — \textquotedblleft{}lake\textquotedblright{}.

\begin{quote}
\emph{gōngyuán lǐ yǒu yí ge hú} — There's a lake in the park.
\end{quote}

Its characters come later; for now it is a word for the ear.

\begin{cousinweb}
\emph{hú} is inherited from Middle Chinese \emph{hu} (Guangyun, in Baxter's transcription), which explains it as a broad body of water.
\end{cousinweb}

\begin{grammarlens}[title={What you've built}]
One more word for describing things.
\end{grammarlens}

\subsection*{Your turn}
\begin{itemize}
\raggedright
  \item \emph{Say it:} \emph{hú}
  \item \emph{Say it:} \emph{hú}, once more
  \item \emph{Say it:} \emph{gōngyuán lǐ yǒu yí ge hú} — There's a lake in the park.
  \item \emph{Recall:} say the Mandarin for to camp out, then the Mandarin for a lawn, then say \emph{hú} again
\end{itemize}

\begin{tcolorbox}[breakable,colback=teal!4,colframe=teal!35!black,title={Before you move on}]
What does hú mean? (\textquotedblleft{}Lake\textquotedblright{}.) Say it once more.
\end{tcolorbox}
\section[sēnlín]{sēnlín — forest}

\label{lesson:ZH-C226-senlin}



\begin{quote}
Before the new one: say the Mandarin for a lawn, then the Mandarin for a lake.
\end{quote}

\subsection*{You'll want to know: sēnlín}
\emph{sēnlín} — \textquotedblleft{}forest\textquotedblright{}.

\begin{quote}
\emph{sēnlín lǐ hěn ānjìng} — It's very quiet in the forest.
\end{quote}

Its characters come later; for now it is a word for the ear.

\begin{cousinweb}
\emph{sēnlín} is built from \emph{sēn} 'dense with trees' and \emph{lín} 'woods'. The second is inherited from Middle Chinese \emph{lim} (Guangyun, in Baxter's transcription), glossed 'trees'; the dictionary describes the first, \emph{srim} there, as how tall trees look.
\end{cousinweb}

\begin{grammarlens}[title={What you've built}]
One more word for describing things.
\end{grammarlens}

\subsection*{Your turn}
\begin{itemize}
\raggedright
  \item \emph{Say it:} \emph{sēnlín}
  \item \emph{Say it:} \emph{sēnlín}, once more
  \item \emph{Say it:} \emph{sēnlín lǐ hěn ānjìng} — It's very quiet in the forest.
  \item \emph{Recall:} say the Mandarin for a lawn, then the Mandarin for a lake, then say \emph{sēnlín} again
\end{itemize}

\begin{tcolorbox}[breakable,colback=teal!4,colframe=teal!35!black,title={Before you move on}]
What does sēnlín mean? (\textquotedblleft{}Forest\textquotedblright{}.) Say it once more.
\end{tcolorbox}
\section[cǎoyuán]{cǎoyuán — grassland; prairie}

\label{lesson:ZH-C226-caoyuan}



\begin{quote}
Before the new one: say the Mandarin for a lake, then the Mandarin for a forest.
\end{quote}

\subsection*{You'll want to know: cǎoyuán}
\emph{cǎoyuán} — \textquotedblleft{}grassland; prairie\textquotedblright{}.

\begin{quote}
\emph{wǒ xiǎng qù cǎoyuán qí mǎ} — I'd love to go riding on the grasslands.
\end{quote}

Its characters come later; for now it is a word for the ear.

\begin{cousinweb}
\emph{cǎoyuán} is built from \zh{草} 'grass' and \emph{yuán} 'a plain', both inherited from Middle Chinese, \emph{tshawX} and \emph{ngjwon} (Guangyun, in Baxter's transcription), which glosses the second 'broad and flat'.
\end{cousinweb}

\begin{grammarlens}[title={What you've built}]
One more word for describing things.
\end{grammarlens}

\subsection*{Your turn}
\begin{itemize}
\raggedright
  \item \emph{Say it:} \emph{cǎoyuán}
  \item \emph{Say it:} \emph{cǎoyuán}, once more
  \item \emph{Say it:} \emph{wǒ xiǎng qù cǎoyuán qí mǎ} — I'd love to go riding on the grasslands.
  \item \emph{Recall:} say the Mandarin for a lake, then the Mandarin for a forest, then say \emph{cǎoyuán} again
\end{itemize}

\begin{tcolorbox}[breakable,colback=teal!4,colframe=teal!35!black,title={Before you move on}]
What does cǎoyuán mean? (\textquotedblleft{}Grassland; prairie\textquotedblright{}.) Say it once more.
\end{tcolorbox}
\section[fēngjǐng]{fēngjǐng — scenery; landscape}

\label{lesson:ZH-C226-fengjing}



\begin{quote}
Before the new one: say the Mandarin for a forest, then the Mandarin for grassland.
\end{quote}

\subsection*{You'll want to know: fēngjǐng}
\emph{fēngjǐng} — \textquotedblleft{}scenery; landscape\textquotedblright{}.

\begin{quote}
\emph{zhèr de fēngjǐng hěn piàoliang} — The scenery here is beautiful.
\end{quote}

Its characters come later; for now it is a word for the ear.

\begin{cousinweb}
\emph{fēngjǐng} is built from \zh{风} 'wind' and \emph{jǐng} 'a view', both inherited from Middle Chinese, \emph{pjuwng} and \emph{kjaengX} (Guangyun, in Baxter's transcription). The dictionary glosses the second 'bright' and 'light'; the sense 'a view' is not given there.
\end{cousinweb}

\begin{grammarlens}[title={What you've built}]
One more word for describing things.
\end{grammarlens}

\subsection*{Your turn}
\begin{itemize}
\raggedright
  \item \emph{Say it:} \emph{fēngjǐng}
  \item \emph{Say it:} \emph{fēngjǐng}, once more
  \item \emph{Say it:} \emph{zhèr de fēngjǐng hěn piàoliang} — The scenery here is beautiful.
  \item \emph{Recall:} say the Mandarin for a forest, then the Mandarin for grassland, then say \emph{fēngjǐng} again
\end{itemize}

\begin{tcolorbox}[breakable,colback=teal!4,colframe=teal!35!black,title={Before you move on}]
What does fēngjǐng mean? (\textquotedblleft{}Scenery; landscape\textquotedblright{}.) Say it once more.
\end{tcolorbox}
